\uuid{HR9T}
\titre{Maximum}
\auteur{bodin ; hanani ; mzari}
\organisation{exo7}
\licence{CC BY-NC-SA 4.0}
\datecreate{2020-01-20}
\niveau{L1}
\module{Analyse}
\chapitre{Propriétés de R}
\sousChapitre{Propriétés des nombres réels}
\difficulte{3}
\theme{Propriétés de nombres réels, maximum}
\exo7qcmid{lille-1.reels.016}

\contenu{
\begin{question}
\qclassification{120.03}
\qtags{motcle=maximum}
Quelles sont les assertions vraies ?
\begin{answers}
    \bad{$\max(x,y) = \frac{x+y-|x|-|y|}{2}$}

    \bad{$\max(x,y) = \frac{x+y-|x+y|}{2}$}

    \bad{$\max(x,y) = \frac{|x+y|-x-y}{2}$}

    \good{$\max(x,y) = \frac{x+y + |x-y|}{2}$}

\end{answers}
\begin{explanations}
On prouve la formule $\max(x,y) = \frac{x+y + |x-y|}{2}$ en distinguant le cas $x-y \ge0$ puis $x-y<0$.
\end{explanations}
\end{question}
}
